现在容易给出隐函数定理对椭圆 PDE 的应用. 在 $\mathbb{R}^3$ 中的有界光滑区域上, 考虑方程
\begin{equation}\label{eq:semilinear-small-data-dirichlet-equation}
\Delta u+f(u)=g(\bm{x})
\end{equation}
并施加 Dirichlet 边界条件. 假设 $f$ 属于 $C^3$ 类, 且 $f(0)=f^{\prime}(0)=0$. 则由 Lemma~\ref{lem:sobolev-pointwise-composition-differentiability}, 映射 $u\mapsto f(u)$ 从 $H^2(\Omega)$ 到自身连续可微. 现在应用逆函数定理可知, 对足够小的 $g\in L^2(\Omega)$, 该方程有解 $u\in H^2(\Omega)\cap H_0^1(\Omega)$. 此外, 该解在范数小的解中唯一. (注意实际上只需要映射 $u\mapsto f(u)$ 作为从 $H^2(\Omega)$ 到 $L^2(\Omega)$ 的映射连续可微. 为此, $f$ 属于 $C^1$ 类就足够.) 如何表述并证明涉及依赖于 $u$ 的导数的非线性和/或非线性边界条件的更一般结果, 是清楚的.

作为分岔问题的一个例子, 考虑方程
\begin{equation}\label{eq:semilinear-dirichlet-eigenvalue-bifurcation-equation}
\Delta u+\lambda u=f(u),
\end{equation}
仍施加 Dirichlet 边界条件, 且 $\Omega$ 和 $f$ 与前面相同. 只要 $\lambda$ 不是 $-\Delta$ 的特征值, 零解就是唯一的小解. 若 $\lambda_0$ 为单特征值, 则前一节的假设适用, 因为 $\Delta+\lambda$ 的值域恰为零空间的正交补; 用 $u_0$ 表示特征函数, 并令 $F(u,\lambda)=\Delta u+\lambda u-f(u)$, 可得 $D_{u\lambda}F(0,\lambda_0)u_0=u_0$. 事实上, Laplacian 的 Dirichlet 问题的第一特征值总是单的.

\begin{Lemma}
\label{lem:first-dirichlet-eigenvalue-simplicity}
设 $\Omega$ 为有界区域. 则 Laplace 方程的 Dirichlet 问题的第一特征值是单的.
\end{Lemma}

\begingroup
\renewcommand{\endProof}{\par\nobreak\hfill$\blacksquare$\par\vspace{0.1cm}}
\begin{Proof}
设 $\lambda_0$ 为最小特征值; 则有
\begin{equation}\label{eq:first-dirichlet-eigenvalue-printed-norm-minimum}
\lambda_0=\min_{\substack{u\in H_0^1(\Omega),\|u\|_2=1}}\|u\|_{1,2}^2.
\end{equation}
回顾 Problem~\ref{ex:sturm-liouville-smallest-eigenvalue-variational-characterization}, 其中对 Sturm-Liouville 问题建立了类似结果. 设 $u_0$ 为极小化元, 并假设 $u_0$ 在 $\Omega$ 内有零点. 现在可以构造分片多项式序列 $p_n$, 它在 $H_0^1(\Omega)$ 中收敛到 $u_0$, 并在 $\Omega$ 的紧子集上一致收敛. 显然有
\begin{equation}\label{eq:first-dirichlet-eigenvalue-piecewise-polynomial-quotients}
\lambda_0=\lim_{n\to\infty}\|p_n\|_{1,2}^2/\|p_n\|_2^2.
\end{equation}
现在 $|p_n|$ 与 $p_n$ 本身具有相同的 $H^1$ 范数和 $L^2$ 范数; 特别地, 当 $n\to\infty$ 时, $|p_n|$ 的 $H^1$ 范数有界. 因而 $|p_n|$ 有一个在 $H_0^1(\Omega)$ 中弱收敛, 且在 $L^2(\Omega)$ 中强收敛的子序列. 由于 $p_n$ 在 $\Omega$ 的紧子集上一致收敛到 $u_0$, 还可得到 $|p_n|\to|u_0|$. 因此有
\begin{equation}\label{eq:first-dirichlet-eigenfunction-absolute-value-quotient}
\bigl\||u_0|\bigr\|_{1,2}^2/\bigl\||u_0|\bigr\|_2^2\leq\lambda_0,
\end{equation}
即 $|u_0|$ 也必为特征函数. 但 $|u_0|$ 非负且在 $\Omega$ 内有零点, 这与最大值原理矛盾.

因此, 关于存在变号特征函数的假设必不成立, 每个特征函数或为正, 或为负. 两个这样的函数不可能彼此正交, 因而该特征值是单的.
\end{Proof}
\endgroup
